\section*{Supplementary Methods}

\subsection*{S1. Hyperparameter Tuning Details}

This section provides comprehensive hyperparameter configurations for all models evaluated in the main text (Methods Section 2.4.3).

\subsubsection*{Linear Models}

\textbf{Linear Regression}: Ordinary least squares implemented via scikit-learn \texttt{LinearRegression} with no hyperparameters to tune (baseline model).

\textbf{Ridge Regression}: L2-regularized linear model with penalty $\alpha\|\mathbf{w}\|_2^2$. Hyperparameter $\alpha$ was tuned via 5-fold cross-validation over the grid $\{0.01, 0.1, 1, 10, 100\}$. Optimal configuration selected based on minimum validation RMSE.

\textbf{Lasso Regression}: L1-regularized linear model with penalty $\alpha\|\mathbf{w}\|_1$ for implicit feature selection. Hyperparameter $\alpha$ was tuned via 5-fold cross-validation over the grid $\{0.001, 0.01, 0.1, 1, 10\}$. Optimal configuration selected based on minimum validation RMSE.

\subsubsection*{Ensemble Methods}

\textbf{Random Forest}: Bagging ensemble of 500 decision trees with randomized feature subsets at each split.

Hyperparameter grid search (5-fold cross-validation):
\begin{itemize}
    \item \texttt{max\_depth}: $\{10, 20, 30, \text{None}\}$ (maximum tree depth)
    \item \texttt{min\_samples\_leaf}: $\{1, 2, 4\}$ (minimum samples per leaf node)
    \item \texttt{max\_features}: $\{\text{sqrt}, \text{log2}, 0.3\}$ (feature sampling strategy)
    \item \texttt{n\_estimators}: 500 (fixed)
    \item \texttt{random\_state}: 42 (fixed)
\end{itemize}

Final configuration: \texttt{max\_depth}=20, \texttt{min\_samples\_leaf}=2, \texttt{max\_features}=sqrt.

\textbf{XGBoost}: Gradient boosted decision trees with L1/L2 regularization.

Hyperparameters optimized via Bayesian optimization (Optuna, 100 trials) with objective function minimizing validation RMSE:
\begin{itemize}
    \item \texttt{learning\_rate}: $[0.01, 0.3]$ (shrinkage applied to each tree)
    \item \texttt{max\_depth}: $[3, 10]$ (maximum tree depth)
    \item \texttt{n\_estimators}: $[100, 1000]$ (number of boosting rounds)
    \item \texttt{subsample}: $[0.6, 1.0]$ (fraction of samples per tree)
    \item \texttt{colsample\_bytree}: $[0.6, 1.0]$ (fraction of features per tree)
    \item \texttt{reg\_alpha}: $[0, 1]$ (L1 regularization weight)
    \item \texttt{reg\_lambda}: $[0, 10]$ (L2 regularization weight)
    \item Early stopping: patience=50 rounds on validation RMSE
    \item \texttt{random\_state}: 42 (fixed)
\end{itemize}

Final configuration: \texttt{learning\_rate}=0.08, \texttt{max\_depth}=6, \texttt{n\_estimators}=450, \texttt{subsample}=0.85, \texttt{colsample\_bytree}=0.7 (default), \texttt{reg\_alpha}=0.1, \texttt{reg\_lambda}=1.5.

\subsubsection*{Support Vector Regression}

\textbf{SVR}: Radial basis function (RBF) kernel with $\epsilon$-insensitive loss.

Hyperparameter grid search (5-fold cross-validation):
\begin{itemize}
    \item $C$ (regularization): $\{0.1, 1, 10, 100\}$
    \item $\epsilon$ (tube width): $\{0.01, 0.1, 0.5\}$
    \item $\gamma$ (RBF kernel coefficient): $\{\text{scale}, \text{auto}, 0.001, 0.01\}$
\end{itemize}

Final configuration selected based on minimum validation RMSE.

\subsubsection*{Neural Network}

\textbf{Multi-layer Perceptron}: Feed-forward neural network with two hidden layers.

Architecture and hyperparameters:
\begin{itemize}
    \item Architecture: [100, 50] (two hidden layers with 100 and 50 neurons)
    \item Activation: ReLU (Rectified Linear Unit)
    \item Dropout rate: 0.2 (applied after each hidden layer for regularization)
    \item Optimizer: Adam with learning rate 0.001
    \item Batch size: 32
    \item Maximum epochs: 500
    \item Early stopping: patience=20 on validation loss
    \item \texttt{random\_state}: 42 (fixed)
\end{itemize}

\subsubsection*{Allometric Baseline}

\textbf{Traditional forestry model}: Power law relationship between LAI and morphological parameters.

Model form: $\text{LAI} = \alpha \times H^\beta \times \text{CD}^\gamma$

Where:
\begin{itemize}
    \item $H$ = tree height (meters)
    \item $\text{CD}$ = crown diameter (meters)
    \item $\alpha, \beta, \gamma$ = parameters estimated via non-linear least squares on training data
\end{itemize}

This model provides an ecological benchmark representing traditional allometric approaches used in forestry applications before machine learning methods.

\subsubsection*{Software Implementation}

All models implemented using:
\begin{itemize}
    \item Python 3.9.7
    \item scikit-learn 1.0.2 (Linear Regression, Ridge, Lasso, Random Forest, SVR, Neural Network)
    \item XGBoost 1.5.0 (gradient boosting)
    \item Optuna 2.10.0 (Bayesian hyperparameter optimization)
    \item NumPy 1.21.2, Pandas 1.3.3 (data manipulation)
\end{itemize}

Random seeds fixed at 42 across all stochastic processes (data splitting, model initialization, cross-validation folds) to ensure reproducibility.

\clearpage
